% ==============================================================================
% Section 4: System Architecture
% ==============================================================================

\section{System Architecture}
\label{sec:architecture}

This section presents the overall architecture of AR-DeC, describing the five principal components and their interactions, the data flow through the system, and the rationale for the selected technology stack.

\subsection{Component Overview}
\label{subsec:components}

AR-DeC is composed of five core modules that collaborate to deliver the integrated learning experience. Each module encapsulates a distinct responsibility while communicating with the others through well-defined interfaces managed by a centralised state layer.

\begin{enumerate}
    \item \textbf{OCR Module}: Captures camera frames, performs on-device text recognition, and outputs structured text with bounding-box coordinates. The module handles image preprocessing (contrast enhancement, skew correction) to maximise recognition accuracy across varying lighting conditions and text formats.

    \item \textbf{ITS Engine}: Receives tokenised words from the OCR module, consults the student model to determine the learner's mastery level for each word, selects an appropriate Bloom's taxonomy level for content delivery, and generates contextual explanations. The engine maintains the domain model (vocabulary database with difficulty ratings and concept relationships) and the tutoring model (content selection algorithms).

    \item \textbf{AR Renderer}: Transforms ITS-generated content into visual overlays registered to the camera feed. The renderer supports multiple presentation modes: text annotations anchored to word positions, floating information cards, and 3D model visualisations for concrete concepts. Rendering strategies are selected based on the ITS engine's interface recommendations.

    \item \textbf{Gamification Engine}: Monitors learning events (scans, reads, quiz responses, daily logins) and applies the points, badges, streaks, and level progression mechanics described in \Cref{sec:implementation}. The engine implements the ARCS-grounded reward schedule and triggers visual feedback (point animations, badge popups) through the AR renderer.

    \item \textbf{Student Model}: Maintains per-learner, per-word mastery estimates using Bayesian Knowledge Tracing. The model persists learning history, tracks vocabulary encounters across sessions, and provides mastery probabilities to the ITS engine for adaptive content selection.
\end{enumerate}

The interactions between these components are illustrated in the system architecture diagram (\Cref{fig:architecture}).

\begin{landscape}
\begin{figure}[htbp]
    \centering
    % ==============================================================================
% TikZ System Architecture Diagram
% ==============================================================================

\begin{tikzpicture}[
    node distance=1.8cm and 2.5cm,
    module/.style={
        rectangle,
        draw=#1!60!black,
        fill=#1!15,
        rounded corners=4pt,
        minimum width=3.2cm,
        minimum height=1.2cm,
        align=center,
        font=\small\bfseries,
        line width=1pt
    },
    input/.style={
        rectangle,
        draw=gray!70,
        fill=gray!10,
        rounded corners=4pt,
        minimum width=2.6cm,
        minimum height=1cm,
        align=center,
        font=\small,
        line width=0.8pt,
        dashed
    },
    arrow/.style={
        -Stealth,
        thick,
        color=#1!70!black
    },
    annot/.style={
        font=\scriptsize\itshape,
        color=#1!60!black
    }
]

% --- Input ---
\node[input] (camera) {Camera Feed};

% --- Core Modules ---
\node[module=blue, right=2cm of camera] (ocr) {OCR Module};
\node[module=purple, right=2cm of ocr] (its) {ITS Engine};
\node[module=teal, right=2cm of its] (ar) {AR Renderer};

% --- Supporting Modules ---
\node[module=orange, below=2cm of its] (gamification) {Gamification\\Engine};
\node[module=red, below=2cm of ocr] (student) {Student Model};

% --- Output ---
\node[input, right=2cm of ar] (display) {AR Display};

% --- Backend ---
\node[input, below=2cm of ar] (firebase) {Firebase\\Backend};

% --- Arrows: Main Pipeline ---
\draw[arrow=gray] (camera) -- node[above, annot=gray] {frames} (ocr);
\draw[arrow=blue] (ocr) -- node[above, annot=blue] {tokens} (its);
\draw[arrow=purple] (its) -- node[above, annot=purple] {content} (ar);
\draw[arrow=teal] (ar) -- node[above, annot=teal] {overlay} (display);

% --- Arrows: Student Model ---
\draw[arrow=red] (student) -- node[left, annot=red] {mastery} (ocr);
\draw[arrow=red] (student) -- node[right, annot=red, xshift=-2mm] {mastery} (its);

% --- Arrows: Gamification ---
\draw[arrow=orange] (its) -- node[right, annot=orange] {events} (gamification);
\draw[arrow=orange] (gamification) -- node[below, annot=orange] {rewards} (ar);

% --- Arrows: Updates ---
\draw[arrow=purple, dashed] (its) -- node[left, annot=purple, xshift=-2mm] {update} (student);
\draw[arrow=orange, dashed] (gamification) -- node[right, annot=orange] {sync} (firebase);
\draw[arrow=red, dashed] (student) -- node[above, annot=red] {persist} (firebase);

% --- Background grouping ---
\begin{scope}[on background layer]
    \node[
        fit=(ocr)(its)(ar)(gamification)(student),
        fill=blue!3,
        draw=blue!20,
        rounded corners=8pt,
        inner sep=12pt,
        label={[font=\footnotesize\scshape, color=blue!40!black]above:AR-DeC Mobile Application}
    ] {};
\end{scope}

\end{tikzpicture}

    \caption{AR-DeC system architecture showing the five core modules and their interactions.}
    \label{fig:architecture}
\end{figure}
\end{landscape}

\subsection{Data Flow}
\label{subsec:data_flow}

The data pipeline follows a sequential-with-branches pattern, as illustrated in \Cref{fig:data_flow}. When a student points their device camera at text, the following sequence executes:

\begin{enumerate}
    \item The camera feed is captured at 30 frames per second and preprocessed for OCR.
    \item The OCR module extracts recognised text blocks along with their spatial coordinates in the camera frame.
    \item A tokeniser segments the extracted text into individual words, filtering stopwords and punctuation.
    \item Each word is looked up in the student model to retrieve the current mastery estimate $P(L_n)$.
    \item Words below the mastery threshold are passed to the ITS engine, which determines the appropriate Bloom's level and generates an explanation.
    \item The AR renderer overlays the explanation content at the word's spatial position in the camera feed.
    \item Simultaneously, the gamification engine records the scan event, awards points, and checks for badge or streak milestones.
    \item The student model is updated with the new interaction data, adjusting mastery probabilities via BKT.
\end{enumerate}

This pipeline is designed to execute within a latency budget of 500 milliseconds from frame capture to overlay display, ensuring a responsive and fluid user experience. The on-device processing of OCR and student model lookups avoids network round-trips for the critical path, with cloud synchronisation occurring asynchronously via Firebase.

\begin{figure}[htbp]
    \centering
    % ==============================================================================
% TikZ Data Flow Diagram
% ==============================================================================

\begin{tikzpicture}[
    node distance=1.2cm,
    step/.style={
        rectangle,
        draw=#1!60!black,
        fill=#1!12,
        rounded corners=3pt,
        minimum width=3cm,
        minimum height=0.9cm,
        align=center,
        font=\small,
        line width=0.8pt
    },
    branch/.style={
        rectangle,
        draw=#1!60!black,
        fill=#1!12,
        rounded corners=3pt,
        minimum width=2.8cm,
        minimum height=0.9cm,
        align=center,
        font=\small,
        line width=0.8pt
    },
    arrow/.style={
        -Stealth,
        thick,
        color=#1!60!black
    },
    note/.style={
        font=\scriptsize\itshape,
        color=gray!70!black
    }
]

% --- Main Pipeline (vertical) ---
\node[step=gray] (camera) {Camera Capture};
\node[step=blue, below=of camera] (ocr) {OCR Processing};
\node[step=blue, below=of ocr] (extract) {Text Extraction};
\node[step=blue, below=of extract] (tokenize) {Word Tokenisation};
\node[step=purple, below=of tokenize] (lookup) {Student Model Lookup};
\node[step=purple, below=of lookup] (its) {ITS Analysis};

% --- Branching outputs ---
\node[branch=teal, below left=1.5cm and 0.5cm of its] (ar_display) {AR Display};
\node[branch=orange, below=1.5cm of its] (gamification) {Gamification\\Update};
\node[branch=red, below right=1.5cm and 0.5cm of its] (student_update) {Student Model\\Update};

% --- Main arrows ---
\draw[arrow=gray] (camera) -- node[right, note] {30 fps} (ocr);
\draw[arrow=blue] (ocr) -- node[right, note] {text blocks + bbox} (extract);
\draw[arrow=blue] (extract) -- node[right, note] {raw strings} (tokenize);
\draw[arrow=blue] (tokenize) -- node[right, note] {filtered tokens} (lookup);
\draw[arrow=purple] (lookup) -- node[right, note] {$P(L_n)$ per word} (its);

% --- Branch arrows ---
\draw[arrow=teal] (its.south) -- ++(0,-0.5) -| (ar_display.north) node[pos=0.25, above, note] {content + position};
\draw[arrow=orange] (its) -- node[right, note] {events} (gamification);
\draw[arrow=red] (its.south) -- ++(0,-0.5) -| (student_update.north) node[pos=0.25, above, note] {interaction data};

% --- Feedback loop ---
\draw[arrow=red, dashed] (student_update.east) -- ++(1.5,0) |- (lookup.east)
    node[pos=0.25, right, note] {updated $P(L_n)$};

% --- Timing annotations ---
\node[note, left=0.3cm of camera] {$t_0$};
\node[note, left=0.3cm of ocr] {$t_0 + 50$ms};
\node[note, left=0.3cm of extract] {$t_0 + 200$ms};
\node[note, left=0.3cm of tokenize] {$t_0 + 220$ms};
\node[note, left=0.3cm of lookup] {$t_0 + 250$ms};
\node[note, left=0.3cm of its] {$t_0 + 350$ms};
\node[note, below=0.1cm of ar_display] {$t_0 + 500$ms};

\end{tikzpicture}

    \caption{Data flow diagram showing the processing pipeline from camera capture to AR display.}
    \label{fig:data_flow}
\end{figure}

\subsection{Technology Stack Justification}
\label{subsec:tech_stack}

The technology choices for AR-DeC were guided by three criteria: cross-platform compatibility, on-device processing capability, and developer ecosystem maturity. \Cref{tab:tech_stack} summarises the selected technologies and their justifications.

\begin{table}[htbp]
\centering
\caption{Technology stack selection and justification.}
\label{tab:tech_stack}
\small
\begin{tabular}{p{3cm} p{3cm} p{6.5cm}}
\toprule
\textbf{Component} & \textbf{Technology} & \textbf{Justification} \\
\midrule
Mobile Framework & React Native 0.83 + Expo 55 & Cross-platform (iOS + Android) from a single TypeScript codebase; extensive ecosystem; over-the-air updates via Expo \\
\addlinespace
AR Engine & ViroReact 2.52 & Native ARKit (iOS) and ARCore (Android) support; React Native integration; 3D model rendering; plane detection for surface anchoring \\
\addlinespace
OCR & react-native-mlkit-ocr (Google ML Kit) & On-device processing (no network dependency); high accuracy across fonts and languages; real-time frame processing; bounding-box output for AR registration \\
\addlinespace
State Management & Zustand 5.0 & Lightweight, performant state with minimal boilerplate; middleware support for persistence; TypeScript-native \\
\addlinespace
Navigation & React Navigation 7 & Industry-standard mobile navigation; native stack animations; bottom tab and modal support \\
\addlinespace
Backend & Firebase (Auth, Firestore, Cloud Functions) & Managed authentication; real-time data synchronisation; serverless cloud functions for AI processing; offline support \\
\addlinespace
AI/NLP & Claude API & High-quality contextual explanations; structured output; ability to adapt explanation level per Bloom's taxonomy \\
\bottomrule
\end{tabular}
\end{table}

The selection of React Native with Expo enables deployment on both iOS and Android from a single codebase, reducing development overhead while reaching the widest possible audience. Google ML Kit's on-device OCR ensures that text recognition operates without network connectivity, which is critical for use in settings with unreliable internet access such as classrooms, libraries, and public transport. ViroReact provides the AR rendering capabilities needed for spatial text overlays while maintaining compatibility with the React Native component model.
